\section{Deletion budgets and graded hook overlap}
\label{rec:sec-hook-budget}

The next method assigns a budget to deleting boxes from a diagram.
Deletion trades a large representation for a smaller diagram inside a
thin hook and a collection of named points.  A graded tensor overlap
estimate controls the hook; falling factorials account for the points.
The budget first minimizes over every subdiagram and clips the result;
the proof then finds a near-minimizer inside a thin hook.

Write $W_\lambda=\log D_\lambda$.  Fix $a=1/4$ and $b=10^{-5}$, and set
\begin{equation}\label{rec:deletion-budget}
 E_n(\lambda)=\min_{\mu\subseteq\lambda}
 \{aW_\mu+b t\log n-t\log(n/t)\},\qquad
 t=n-|\mu|,\qquad E_n^+=\max(0,E_n).
\end{equation}
The $t$-dependent logarithmic terms have continuous value zero at
$t=0$.  All traces below
are ordinary, unnormalized traces.

\begin{theorem}\label{rec:deletion-moment}
There are real exponents $p(d)\ge1$, indexed by integers $d\ge1$,
with $\sup_{d\ge1}p(d)<\infty$ such that, for every integer $d\ge1$
and every partition $\lambda\vdash n=2^d$,
\[
 D_\lambda\Tr|B_n|^{2p(d)}\le e^{E_n^+(\lambda)}.
\]
There is a fixed integer $d_*\ge1$ above which one may take
\[
 p(d)=(1+d^{-1/2})
       \max\{p(\lfloor d/2\rfloor),p(\lceil d/2\rceil)\}.
\]
In particular $E_n^+\le aW_\lambda$ yields a fixed negative power
of $D_\lambda$ for the sweep norm.
\end{theorem}

\subsection{Whitening a graded tensor overlap}

Consider an $a_0$ by $b_0$ grid with $n=a_0b_0$ and $t$ deleted
cells.  Put $s=n-t$.  When $s=0$, the surviving representation is
one-dimensional and the overlap bound below is immediate; thus the
whitening argument assumes $s>0$.  Suppose $\mu\vdash s$ lies in an $(r,r)$ hook,
where $1\le r\le n^{1/16}$ is an integer. Let $X,Y$ be products of
positive row and column group-algebra elements, supported on the
local type tuples $\gamma,\delta$, respectively. Then
\begin{equation}\label{rec:graded-hook-overlap}
 \Tr_{V_\mu}XY\le\Tr_{V_\gamma}X\,\Tr_{V_\delta}Y\,
 \exp\{O((a_0+b_0+1)r^2\log(n+1)+t)+\min(0,W_\gamma+W_\delta-W_\mu)\}.
\end{equation}
Here local dimensions and traces multiply over the lines, so
$W_\gamma=\sum_i\log D_{\gamma_i}$ and
$W_\delta=\sum_j\log D_{\delta_j}$. For the
balanced grids used in the induction, $a_0+b_0=O(\sqrt n)$, the
error is $O(n^{4/5}+t)$. The displayed formula retains the error for
arbitrary rectangular geometry.

We supply the entropy gain in this inequality.  Realize the hook
inside signed tensors over $\mathbb C^r\oplus\mathbb C^r$, with the
second summand odd.  A compatible compact type has even and odd
weight lists $\xi,\eta$.  Its entropy $h$ differs from the symmetric
group log dimension by $O(r^2\log n)$; carrier dimensions and the
number of compatible types cost $\exp O(r^2\log n)$.  Over all grid
lines these errors are $O((a_0+b_0+1)r^2\log(n+1))$. The complete
positive-factor estimate, including all carrier multiplicities, is
Proposition~\ref{s:one-sided-marked-overlap}, applied with $q=r$.
Its local hook hypothesis causes no additional restriction here:
a local type that occurs in the restriction of $V_\mu$ is a
subdiagram of $\mu$ by the Littlewood--Richardson rule, and hence
is itself an $(r,r)$-hook shape. Any other local type acts as zero
on $V_\mu$. Thus the proposition proves
\eqref{rec:graded-hook-overlap}. We retain the following deformation
calculation to explain the entropy gain; the proposition supplies
the passage to arbitrary positive factors and their multiplicities.

For a local type with list of weights $(p_i)$ and $u$ sites, use the
strictly positive density matrix with spectrum
$(p_i+1)/(u+2r)$.  Twirl it over the even and odd unitary groups.
The highest-weight term in the twirl dominates the local type
projection with loss $e^{W_\gamma+O((a_0+b_0+1)r^2\log(n+1))}$, and similarly in the
other direction.  This regularization permits inverse square roots
even when some weights vanish.

Work in one compatible global compact sector at a time.  For each
pair $(\xi,\eta)$ with at most $r$ parts in each partition and
$|\xi|+|\eta|=s$ for which $V_\mu$ occurs in
$\operatorname{Ind}_{S_{|\xi|}\times S_{|\eta|}}^{S_s}
(V_\xi\otimes V_{\eta^{\mathsf t}})$, put
\[
 \begin{gathered}
 K_r=U(r)_{\rm even}\times U(r)_{\rm odd},\\
 \mathcal U_\tau=
 \mathbf S_\xi(\mathbb C^r)\otimes\mathbf S_\eta(\mathbb C^r),\\
 \tau=\tau_{(\xi,\eta)}:K_r\longrightarrow U(\mathcal U_\tau).
 \end{gathered}
\]
Here $\mathbf S_\xi$ and $\mathbf S_\eta$ denote the irreducible
polynomial carrier representations.  The action $\tau$ extends
polynomially to the invertible parity-preserving block matrices.
Write $\chi_\tau(B)=\Tr_{\mathcal U_\tau}\tau(B)$ for its character,
and
\[
 h=s\log s-\sum_i\xi_i\log\xi_i-\sum_j\eta_j\log\eta_j,
\]
with zero summands omitted.  On the whole signed tensor space let
$\rho_s(B)=B^{\otimes s}$ be the color action.  It is distinct from
the carrier action $\tau(B)$; the position action remains the signed
permutation action.  Let $Z=Z_\tau$ be the full compact
$\tau$-isotypic projection in this tensor space, including every
multiplicity copy.

Let $\mathcal W$ be the set of $s$ occupied cells.  For the row and
column densities from the local twirls, write
\[
 \begin{gathered}
 R_{\mathcal W}=\bigotimes_{(i,j)\in\mathcal W}\rho_i,\\
 C_{\mathcal W}=\bigotimes_{(i,j)\in\mathcal W}\sigma_j,\\
 \bar\rho=a_0^{-1}\sum_{i=1}^{a_0}\rho_i,\qquad
 A=\bar\rho^{-1/2}.
 \end{gathered}
\]
An empty line may be assigned any strictly positive parity-preserving
density of trace one.  The regularized densities make $\bar\rho$
strictly positive, so $A$ is an invertible parity-preserving block
matrix.  With $dg$ denoting normalized Haar measure on $K_r$, the
deformed character formula is
\begin{equation}\label{rec:graded-deformed-projector}
 Z=(\dim\tau)\int_{K_r}
          \chi_\tau((Ag)^{-1})\rho_s(Ag)\,dg.
\end{equation}
Apply Lemma~\ref{s:compact-coefficient-extraction} with carrier matrix
$\tau(A)^{-1}$. Its coefficient integral acts as
$\tau(A)^{-1}$ on the carrier of every copy of $\tau$ and as zero
on the other compact types. Multiplication on the left by
$\rho_s(A)$ therefore gives $Z$, proving the displayed formula.
For unitary $g$, the character factor is
\[
 \begin{gathered}
 \chi_\tau((Ag)^{-1})
 =\Tr_{\mathcal U_\tau}\bigl(\tau(g)^*\tau(A)^{-1}\bigr),\\
 |\chi_\tau((Ag)^{-1})|\le(\dim\tau)e^{-h/2}.
 \end{gathered}
\]
The bound follows by evaluating the highest-weight monomials of
$A^{-1}$ and using the entropy maximizing proportions.  The factor
$\dim\tau$ outside the integral is separate from the one in this
character bound.

For each cell define the squared Hilbert--Schmidt quantity
\[
 \begin{aligned}
 z_{ij}(g)&=\|\rho_i^{1/2}Ag\sigma_j^{1/2}\|_{\HS}^{\,2}\\
          &=\Tr(A\rho_iA\,g\sigma_jg^*)\ge0.
 \end{aligned}
\]
Since $A\bar\rho A=I$ and every $\sigma_j$ has trace one, its sum over
the full grid is
\[
 \begin{aligned}
 \sum_{i=1}^{a_0}\sum_{j=1}^{b_0}z_{ij}(g)
 &=a_0\sum_{j=1}^{b_0}\Tr(g\sigma_jg^*)\\
 &=a_0b_0=n.
 \end{aligned}
\]
Arithmetic--geometric means on the occupied cells therefore give
\[
 \begin{aligned}
 \left\|\bigotimes_{(i,j)\in\mathcal W}
       \rho_i^{1/2}Ag\sigma_j^{1/2}\right\|_{\HS}^{\,2}
 &=\prod_{(i,j)\in\mathcal W}z_{ij}(g)\\
 &\le(n/s)^s\le e^t.
 \end{aligned}
\]
Thus the tensor coefficient has Hilbert--Schmidt norm at most
$e^{t/2}$.  The triangle inequality in
\eqref{rec:graded-deformed-projector}, together with the character
bound, gives the exact scale
\[
 \|R_{\mathcal W}^{1/2}ZC_{\mathcal W}^{1/2}\|_{\HS}
 \le(\dim\tau)^2e^{-h/2+t/2}.
\]
Its square is at most $(\dim\tau)^4e^{-h+t}$.  The two local twirl
domination coefficients enter the overlap norm through their square
roots.  After squaring, their logarithms add
$W_\gamma+W_\delta+O((a_0+b_0+1)r^2\log(n+1))$ to this squared
logarithmic bound.  Applying this estimate to each compatible global
type, retaining all carrier dimensions and the number of such types,
and using $h=W_\mu+O(r^2\log(n+1))$ gives the entropy exponent
$W_\gamma+W_\delta-W_\mu+O((a_0+b_0+1)r^2\log(n+1)+t)$.

The other branch of the minimum also needs the carrier factors.
On the signed tensor space, let $P_\mu$ be the full symmetric-group
isotypic projection and let $w_\mu\ge1$ be its multiplicity.
Since $P_\mu$ commutes with $X$ and $Y$,
\[
 w_\mu\Tr_{V_\mu}(XY)
 =\Tr(P_\mu XY)\le(\Tr X)(\Tr Y).
\]
Each trace on the right is a product of local signed-tensor traces.
A local trace is its one-copy irreducible trace multiplied by the
local carrier multiplicity, which is at most
$(n+r+1)^{10r^2}$ by Lemma~\ref{lem:signed-hook-carriers}.
Keeping these multiplicities converts the global signed-tensor
traces to the local irreducible traces in
\eqref{rec:graded-hook-overlap}, and gives its zero entropy branch
with the same allowed error.

\subsection{The sparse bound and near-optimal deletion}

For $\lambda=(n-k,\beta)$, a separate path argument gives
\begin{equation}\label{rec:hook-forest-sparse}
 \|B_n^*B_n|_{V_\lambda}\|_{\op}
 \le e^{Ck}(2kd/n)^{k/2},\qquad 2kd/n\le1.
\end{equation}
Apply the detailed weighted-forest proof of \eqref{ten:spin:eq:5},
with $\Delta=2kd/n$. Its inclusion-probability argument and weighted
neighbor bound hold for every $k\le n/(2d)$, not only for a fixed
polynomial sparse range. In that proof the sum over forests is
\[
 e^{Ck}\sum_{1\le j\le k/2}\frac{k^j}{j!}\Delta^{k-j}.
\]
When $\Delta\le1$, this is at most
$e^{Ck}\Delta^{k/2}\sum_{j\ge0}k^j/j!
\le e^{(C+1)k}\Delta^{k/2}$.
Thus the same row-sum estimate gives
$\|B_n\|^2\le e^{C'k}\Delta^{k/2}$ throughout the full displayed
range. Since $\|B_n^*B_n\|=\|B_n\|^2$, this proves
\eqref{rec:hook-forest-sparse}. For $k=1$ the centered kernel is zero,
as its one-vertex encounter graph is isolated; for $k=0$ the claimed
sparse estimate is unnecessary.

The budget has the elementary bounds
\begin{equation}\label{rec:deletion-basic-bounds}
 E_n^+\le aW_\lambda,\qquad
 0\le E_n^+-E_n\le Cn^{1-b}.
\end{equation}
The second follows by minimizing $bt\log n-t\log(n/t)$ over
$0\le t\le n$.  A near-minimizing diagram can be chosen in the
$(r,r)$ hook with $r=\lfloor n^{1/16}\rfloor$, at additional cost
$O(n^{99/100})$, and its deletion size obeys
\begin{equation}\label{rec:deletion-size-control}
 t\le n^{1-b/2}+C(W_\lambda+n^{99/100})/\log n.
\end{equation}
To prove this carefully, set $g_n(t)=bt\log n-t\log(n/t)$,
with $g_n(0)=0$. For $t>0$,
\[
 g_n'(t)=b\log n+\log(t/n)+1\le b\log n+1.
\]
Choose an exact minimizer $\mu_0\subseteq\lambda$, and let $v_j$
be the number of its boxes below row $j$ and right of column $j$.
They form a translated straight partition $\nu_j$. Let $\mu^{(j)}$
be the diagram left after deleting that core. Filling $\mu^{(j)}$
first and then $\nu_j$ gives
$D_{\mu_0}\ge D_{\mu^{(j)}}D_{\nu_j}$.
Among the $v_j$ core boxes, the $v_{2j}$ boxes beyond its first $j$
rows and columns have hooks at most $2v_j/j$; all other hooks are at
most $v_j$. Thus the hook formula and $v!\ge(v/e)^v$ yield
\[
 \log D_{\nu_j}\ge v_{2j}\log(j/2)-v_j.
\]
If $v_{2j}\ge v_j/2>0$, deleting this core changes the objective by
at most
\[
 -a\{(v_j/2)\log(j/2)-v_j\}+v_j(b\log n+1).
\]
For $j\ge\lceil n^{1/32}\rceil$ this is negative at all sufficiently
large $n$, since $b-a/64<0$. Exact minimality therefore forces
$v_{2j}<v_j/2$ whenever $v_j>0$. Starting at
$j_0=\lceil n^{1/32}\rceil$ and iterating dyadically to the largest
$j_h\le r$ gives $j_h>r/2$ and
\[
 v_r\le v_{j_h}\le n2^{-h}\le2nj_0/r=O(n^{31/32}).
\]
Truncate $\mu_0$ once at $r$. Its dimension decreases, while the
derivative bound makes the increase of $g_n$ at most
$O(n^{31/32}\log n)=O(n^{99/100})$.
For its deletion size $t\ge n^{1-b/2}$ one has
$g_n(t)\ge(b/2)t\log n$. The new objective is at most
$E_n(\lambda)+O(n^{99/100})\le aW_\lambda+O(n^{99/100})$,
and its dimension term is nonnegative. This proves
\eqref{rec:deletion-size-control}. Finally $g_n$ has minimum
$-n^{1-b}/e$, establishing the second bound in
\eqref{rec:deletion-basic-bounds} even when $E_n$ is negative.

\subsection{The cancellation of named-point factors}

Induce from the selected hook diagram $\mu$ on the remaining $n-t$
points, fixing each of $t$ distinct named points individually.  The
multiplicity of $V_\lambda$ in this induced module is
$f^{\lambda/\mu}$ by branching.  Splitting tableau entries also gives
\begin{equation}\label{rec:hook-induction-dimension}
 D_\lambda\le\binom nt D_\mu f^{\lambda/\mu}.
\end{equation}
We specify precisely the coefficient restriction used on a row
group.  Let $G=\prod_i S_{l_i}$, let $\alpha$ be a tuple of row
types, and let $X$ be a positive group-algebra element whose only
nonzero Fourier block is $X_\alpha\ge0$.  Use the conventions
\[
 \widehat x(\alpha)=\sum_{g\in G}x(g)\rho_\alpha(g),\qquad
 x(g)=\frac{D_\alpha}{|G|}
               \Tr\{X_\alpha\rho_\alpha(g^{-1})\}.
\]
For a fixed set $S$ of named positions, with
$s_i=|S\cap\operatorname{row}_i|$, let
$H=\{g\in G:g(s)=s\text{ for every }s\in S\}$ be its pointwise
stabilizer.  Thus $H=\prod_i S_{l_i-s_i}$.  Define
\begin{equation}\label{rec:stabilizer-expectation}
 X_S=E_HX=\sum_{h\in H}x(h)h,\qquad
 X_{S,\gamma}=\sum_{h\in H}x(h)\rho_\gamma(h).
\end{equation}
This keeps the original coefficients on the subgroup $H$; it
neither prescribes arbitrary ordered images nor averages coefficients
over other cosets.

The restriction is positive.  Indeed the matrix of regular convolution
by $X$ has entries $x(g_1g_2^{-1})$.  Its compression to the basis
indexed by $H$ is exactly regular convolution by $E_HX$, and a
compression of a positive operator is positive.  Equivalently,
write the subgroup restriction as
$V_\alpha|_H=\bigoplus_\gamma V_\gamma\otimes M_{\alpha\gamma}$
and let $\Pi_\gamma$ be the corresponding projection.  Matrix
coefficient orthogonality gives the exact formula
\[
 X_{S,\gamma}=
 \frac{D_\alpha|H|}{|G|D_\gamma}
 \Tr_{M_{\alpha\gamma}}
                  (\Pi_\gamma X_\alpha\Pi_\gamma),
\]
which is positive and is zero unless every
$\gamma_i\subseteq\alpha_i$.  Fourier inversion at the identity,
or summing this partial-trace formula, gives
\begin{equation}\label{rec:stabilizer-trace-identity}
 \sum_\gamma D_\gamma\Tr X_{S,\gamma}
 =|H|x(e)
 =\frac{|H|}{|G|}D_\alpha\Tr X_\alpha.
\end{equation}
Each summand is nonnegative.  Since
$|G|/|H|=\prod_i(l_i)_{s_i}$, it follows that
\begin{equation}\label{rec:coefficient-restriction}
 D_\gamma\Tr X_{S,\gamma}\le
 \frac{D_\alpha\Tr X_\alpha}{\prod_i(l_i)_{s_i}}.
\end{equation}
For factorizable $X$, coefficient restriction factors over the rows.
The column operator $Y$ has the identical construction with column
counts $s'_j$ and its own pointwise stabilizer.

These are exactly the restrictions required by the induced trace.
Its diagonal coset blocks are indexed by ordered positions of the
$t$ named points.  A product of a row and a column permutation, in
either order, can return a named point $(i,j)$ to itself only if both
fix it: the row permutation changes only the column coordinate, and
the column permutation changes only the row coordinate.  Neither can
undo a coordinate change made by the other, so the intermediate
position must also be $(i,j)$.  This
holds for every named point simultaneously.  Consequently, in the
diagonal block associated with the position set $S$, only the
pointwise-stabilizing coefficients of both operators occur.
Identifying the other positions with the carrier of $\mu$ gives
that block's trace as $\Tr_{V_\mu}(X_SY_S)$.  The choice of coset
representative merely conjugates this identification.  Thus
\begin{equation}\label{rec:induced-diagonal-trace}
 \Tr_{\operatorname{Ind}V_\mu}(XY)
   =t!\sum_{|S|=t}\Tr_{V_\mu}(X_SY_S).
\end{equation}
All these traces are nonnegative because $X_S,Y_S\ge0$.
Likewise the irreducible traces of $XY$ are nonnegative.  The
branching multiplicity therefore implies
$f^{\lambda/\mu}\Tr_{V_\lambda}(XY)
\le\Tr_{\operatorname{Ind}V_\mu}(XY)$.
Together with \eqref{rec:hook-induction-dimension}, this supplies
the factor $D_\mu\binom nt t!$ used below.  In particular, the
argument uses diagonal induced blocks and positivity of the
pointwise-stabilizer expectation, rather than positivity of an
arbitrary off-diagonal coset restriction.

For each child type $\alpha$ and every prescribed stabilizer
subshape $\gamma\subseteq\alpha$ with $h$ deleted positions, the
minimum defining the child budget and its clipping bound give
\[
 E_m^+(\alpha)\le aW_\gamma+bh\log m-h\log(m/h)+Cm^{1-b}.
\]
The candidate $\gamma$ need not minimize the child objective.
Summing over the two balanced axes gives a total clipping error
$O(n^{1-b/2})$. Consequently the sum of child budgets is at most
\begin{equation}\label{rec:hook-local-budget}
 aU+bt\log n-J+O(n^{1-b/2}),
 \qquad U=W_\gamma+W_\delta,
\end{equation}
where the count entropy is
\[
 J=\sum_i s_i\log(b_0/s_i)+\sum_j s'_j\log(a_0/s'_j).
\]
Each margin sums to $t$, and zero summands are omitted.  The number of compatible marked-cell sets with these
margins is at most
\begin{equation}\label{rec:hook-margin-count}
 \frac{t!}{\prod_i s_i!\prod_j s'_j!}
 \le\exp\{J-t\log(n/t)+O(t)\}.
\end{equation}
This is the usual assignment count: assign the labels to row slots
and column slots, then forget their order within each line.  Stirling
gives the second inequality.  The factor $\binom nt t!$ from
\eqref{rec:hook-induction-dimension} cancels the falling factorials
in \eqref{rec:coefficient-restriction}, up to $e^{O(t)}$.
More precisely, $(m)_h\ge(m/e)^h$ and both margins sum to $t$, so
\[
 \frac{\binom nt t!}
 {\prod_i(b_0)_{s_i}\prod_j(a_0)_{s'_j}}
 \le e^{2t}.
\]
The cancellation must precede the estimate of the number of profiles;
otherwise it would leave a spurious factorial cost.

Apply \eqref{rec:graded-hook-overlap}.  Its raw trace bound has gain
$\min(0,U-W_\mu)$.  The child quantities being propagated are
dimension-weighted traces.  Dividing their dimensions out and
restoring the parent dimension multiplies this raw bound by
$e^{W_\mu-U}$.  The resulting exponent is
$W_\mu-U+\min(0,U-W_\mu)=\min(0,W_\mu-U)$.
It combines with \eqref{rec:hook-local-budget} by the elementary inequality
\[
 aU+\min(0,W_\mu-U)\le aW_\mu\qquad(0<a<1).
\]
Together with \eqref{rec:hook-margin-count}, this recovers exactly
the expression minimized in \eqref{rec:deletion-budget}, with errors
$O(n^{1-b/2}+t+n^{99/100})$. The logarithm of the number of remaining
types and profiles is $O(n^{3/4}+\sqrt n\log(n+1))$: the full child
partitions and stabilizer subpartitions each contribute
$O(a_0\sqrt{b_0}+b_0\sqrt{a_0})$ by the partition bound, and the
row and column margins contribute $O((a_0+b_0)\log(n+1))$.
These terms are included in the stated errors.

\begin{proof}[Proof of Theorem~\ref{rec:deletion-moment}]
Write $Q_n=B_n^*B_n$ and, on $V_\lambda$, put
\[
 \begin{gathered}
 T_n^\lambda(P)=D_\lambda\Tr Q_n^P
                =D_\lambda\Tr|B_n|^{2P}\quad(P>0),
 \\
 \theta_{d,\lambda}=\|Q_n|_{V_\lambda}\|_{\op}
       =\|B_n|_{V_\lambda}\|_{\op}^2.
 \end{gathered}
\]
The powers in $T_n^\lambda(P)$ are positive real spectral powers.
The eigenvalues of $Q_n$
lie in $[0,1]$. Thus, when $\theta_{d,\lambda}>0$ and $v\ge u>0$,
\begin{equation}\label{rec:deletion-spectral-slack}
 \begin{gathered}
 T_n^\lambda(v)\le\theta_{d,\lambda}^{v-u}T_n^\lambda(u),
 \\
 \theta_{d,\lambda}^u\le e^{-W_\lambda}T_n^\lambda(u),
 \\
 T_n^\lambda(u)\le e^{2W_\lambda}\theta_{d,\lambda}^u.
 \end{gathered}
\end{equation}
The last inequality includes one factor $D_\lambda$ for the number
of eigenvalues in the unnormalized trace and one for the weight in
$T_n^\lambda$.

The trivial shape has $D_\lambda=1$ and $B_n=I$. Its undeleted
budget choice has value zero, so $E_n^+=0$ and its target is the
equality $1=1$. Every block with $B_n(\lambda)=0$ has zero moment
for $P>0$. In particular this handles all shapes in
Lemma~\ref{lem:annihilation}, without requiring a converse to that
lemma. Every remaining block is nontrivial and has
$0<\theta_{d,\lambda}<1$ by Lemma~\ref{lem:finitegap}.

We first fix the exponent requirement for the sparse range. Let
$C_f\ge0$ be the absolute constant in
\eqref{rec:hook-forest-sparse}. For all sufficiently large $d$,
\[
 2d\le n^{b/16},\qquad C_f\le(b/64)\log n.
\]
For $1\le k\le n^{1-b/8}$ these inequalities give
$2kd/n\le n^{-b/16}\le1$ and hence
$\theta_{d,\lambda}\le n^{-bk/64}$. Since
$D_\lambda\le n^k$, the last inequality in
\eqref{rec:deletion-spectral-slack} gives
\[
 T_n^\lambda(P)\le n^{(2-bP/64)k}\le1\le e^{E_n^+(\lambda)}
 \qquad\text{when }P\ge P_{\rm sp}:=128/b.
\]
These size thresholds are independent of $P$.

For every remaining nontrivial, nonzero shape with $k>n^{1-b/8}$,
the first column has length at most $n/2$ by
Lemma~\ref{lem:annihilation}. The dimension estimates at the start of
Section~\ref{ten:spin} therefore give
$W_\lambda\ge c n^{1-b/8}$.

We now apply the preceding static calculation to a sweep moment.
Split the coordinates into batches of $\lfloor d/2\rfloor$ and
$\lceil d/2\rceil$ bits, giving a grid with
$a_0=2^{\lfloor d/2\rfloor}$ rows and
$b_0=2^{\lceil d/2\rceil}$ columns. Write
$B_n=B_HB_V$, where $B_H$ and $B_V$ are products of the row and
column sub-sweeps, respectively. At a common real child exponent
$q\ge1$, put
\[
 X=(B_H^*B_H)^q,\qquad Y=(B_VB_V^*)^q.
\]
These are positive group-algebra elements, each a product over its
lines. The spectral comparison and Araki--Lieb--Thirring inequality
used in \eqref{ten:spin:eq:7} apply for every real $q\ge1$ and give
\[
 T_n^\lambda(q)\le D_\lambda\Tr_{V_\lambda}(XY).
\]
Decompose $X$ and $Y$ into their full local irreducible blocks.
For a row profile $\alpha$, the dimension-weighted trace of its
$i$th block is $T_{b_0}^{\alpha_i}(q)$, and for a column profile
$\beta$ it is $T_{a_0}^{\beta_j}(q)$. The two orientations of a
child's positive square have the same trace. Increasing a child
exponent decreases its singular-power trace, so the child targets
at any assigned exponents no larger than $q$ bound these quantities.

For each pair of profiles, coefficient restriction to the pointwise
stabilizer of a marked set gives the positive factors considered
above. Decomposing these restricted factors into their local types
and applying \eqref{rec:graded-hook-overlap} gives exactly the
dimension factor and child-budget comparison following
\eqref{rec:hook-margin-count}. Summing the profiles and marked sets
therefore yields
\[
 \log T_n^\lambda(q)
 \le E_n^+(\lambda)
       +C\bigl(n^{1-b/2}+t+n^{99/100}\bigr).
\]
The deletion-size estimate and the lower bound for $W_\lambda$
make this error small relative to the dimension:
\[
 \begin{split}
 n^{1-b/2}+t+n^{99/100}
 &\le C\left(n^{1-b/2}+n^{99/100}
                  +\frac{W_\lambda}{\log n}\right)\\
 &\le C'\frac{W_\lambda}{d}.
 \end{split}
\]
Indeed $\log n=d\log2$, and the fixed gaps between $1-b/8$ and
both $1-b/2$ and $99/100$ absorb the additional factor $d$ for all
sufficiently large $d$. Consequently there is an absolute
$C_r\ge0$ such that
\begin{equation}\label{rec:deletion-preliminary-moment}
 T_n^\lambda(q)\le
 \exp\{E_n^+(\lambda)+C_rW_\lambda/d\}
\end{equation}
for all sufficiently large $d$. Its constant is independent of $q$.

Choose an integer $d_*\ge1$ beyond these thresholds and the sparse
thresholds, and large enough that, for $d>d_*$,
\[
 C_r/d\le1/4,\qquad C_r/d\le1/(2\sqrt d).
\]
For example the latter two requirements hold once $d_*$ is at least
$\max\{\lceil4C_r\rceil,\lceil4C_r^2\rceil\}$.
Since $E_n^+\le W_\lambda/4$, the middle inequality in
\eqref{rec:deletion-spectral-slack} and
\eqref{rec:deletion-preliminary-moment} imply
\[
 \theta_{d,\lambda}^q
 \le\exp\{-(3/4-C_r/d)W_\lambda\}
 \le e^{-W_\lambda/2}.
\]
The first inequality in \eqref{rec:deletion-spectral-slack} now shows
\begin{equation}\label{rec:deletion-error-payment}
 T_n^\lambda(q(1+d^{-1/2}))
 \le\exp\{E_n^+(\lambda)
        +(C_r/d-1/(2\sqrt d))W_\lambda\}
 \le e^{E_n^+(\lambda)}.
\end{equation}
Thus the increase in the stated recursion pays the entire preliminary
error, with all thresholds chosen before the finite exponent below.

It remains to start this induction for every partition. Define the
finite set
\[
 \mathcal F_*=
 \{(d,\lambda):1\le d\le d_*,\ \lambda\vdash2^d,
    \ \lambda\ne(2^d),\ \theta_{d,\lambda}>0\}.
\]
The norm-gap lemma makes every $-\log\theta_{d,\lambda}$ on this
set strictly positive. With an empty inner maximum interpreted as
zero, choose the single real number
\begin{equation}\label{rec:deletion-finite-base}
 P_*=\max\left\{1,P_{\rm sp},
  \max_{(d,\lambda)\in\mathcal F_*}
   \frac{2W_\lambda-E_{2^d}^+(\lambda)}
        {-\log\theta_{d,\lambda}}\right\}.
\end{equation}
This number is finite, since $\mathcal F_*$ is finite; its numerators
are nonnegative because $E_n^+\le W_\lambda/4$.
For every member of $\mathcal F_*$ and every $P\ge P_*$, the last
inequality in \eqref{rec:deletion-spectral-slack} gives
\[
 T_n^\lambda(P)
 \le\exp\{2W_\lambda-P(-\log\theta_{d,\lambda})\}
 \le e^{E_n^+(\lambda)}.
\]
The trivial and zero blocks were already handled directly. Thus this
one choice proves every target for $1\le d\le d_*$, including all
partitions at each starting size.

Set $p(d)=P_*$ for $1\le d\le d_*$ and use the displayed recursion
in the theorem for $d>d_*$. Both child indices are positive and
smaller whenever the recursion is used. Strong induction, using the
sparse estimate or \eqref{rec:deletion-error-payment} at the larger child exponent, proves the target
for every $d\ge1$ and every partition.

To check boundedness with the rounded child indices, follow a child
attaining the maximum until the path $d^{(0)},\ldots,d^{(h)}$ first
reaches $d^{(h)}\le d_*$. If $h>0$, then
$d^{(h-1)}\ge d_*+1$ and
$d^{(j)}-1\ge2(d^{(j+1)}-1)$. Hence
\[
 \sum_{j=0}^{h-1}(d^{(j)})^{-1/2}
 \le d_*^{-1/2}\sum_{r\ge0}2^{-r/2}
 =\frac{2+\sqrt2}{\sqrt{d_*}}.
\]
It follows that
\[
 p(d)\le P_*\exp\left(\frac{2+\sqrt2}{\sqrt{d_*}}\right)
 \qquad(d\ge1).
\]
Thus the real exponents form a bounded sequence.

For completeness, take an integer $R\ge6\sup_dp(d)$ for the number
of forward sweeps. The target implies
$\theta_{d,\lambda}^{p(d)}\le D_\lambda^{-3/4}$, and
\eqref{rec:deletion-spectral-slack} then gives, on each nonzero block,
\[
 T_n^\lambda(R)
 \le D_\lambda^{1-(3/4)R/p(d)}
 \le D_\lambda^{-7/2}\le D_\lambda^{-2}.
\]
Schatten H\"older for the $R$ forward factors bounds
$\|B_n^R\|_{\HS}^2$ by $\Tr Q_n^R$. The regular Fourier sum over
surviving nontrivial types therefore tends to zero by
Lemma~\ref{lem:degrees}, as asserted. This step requires $R$ to be
an integer, but imposes no integrality condition on $p(d)$.
\end{proof}
